\section{本章小结}
\label{sec:gcc-summary}

本章回答了章首的问题：独立高斯噪声下，每次使用在功率预算$P$内最多能可靠送出多少比特。答案由三层论证组成。

第一层是上界。输入二阶矩有限、噪声独立且方差为正时，输出有密度，互信息等于输出熵减去噪声熵；输出方差不超过$P+\sigma^2$，固定方差时高斯分布的熵最大，于是单次互信息不超过$\frac12\log_2(1+P/\sigma^2)$。去掉任何一个条件，结论都可能失效：不限功率或噪声方差为零时，容量没有上限。

第二层把上界连接到可靠传输。零均值高斯输入$\mathcal{N}(0,P)$使输出仍为高斯并达到上界，它是唯一的最优输入。坐标独立服从$\mathcal{N}(0,P')$、$P'<P$的随机码配合信息密度的阈值判决，使发送的码字以接近一的概率越过门限、其他码字以指数小的概率越过门限；超预算和错误偏高的码字在删去一半消息时一并去掉，得到满足逐条功率约束、最大错误趋于零的码。反方向由Fano不等式、逐时刻的熵上界和容量对功率的凹性给出。两者合起来，$C_{\mathrm{op}}(P)=\frac12\log_2(1+P/\sigma^2)$比特每次实值使用，单位换成复值使用或每秒时须一起换算功率和使用次数。

第三层处理多个分量。独立的平行分量把单次问题化为功率分配，最优分配是注水：低噪声分量先得到功率，所有活跃分量的边际收益相等，噪声高于水位的分量不用。它与率失真中的反注水形式相似，但分配的是发送功率而不是允许的失真。相关噪声经正交变换变成独立分量，欧氏预算保持不变，容量是对噪声协方差的特征值注水；白化同样可用，但必须把预算一起变换成带权的形式。

\begin{BookInlineTable}
 \centering
 \caption{本章的三种高斯信道}
 \label{tab:gcc-summary}
 \begin{tabular}{llll}
  \toprule
  模型 & 噪声 & 容量 & 最优输入 \\
  \midrule
  标量 & $\mathcal{N}(0,\sigma^2)$ & $\frac12\log_2(1+P/\sigma^2)$ & $\mathcal{N}(0,P)$ \\
  平行 & 独立，方差$\sigma_j^2$ & 对$\sigma_j^2$注水 & 独立，功率$(\nu-\sigma_j^2)^{+}$ \\
  相关 & $\mathcal{N}(\bm{0},\bm{\Sigma})$ & 对特征值$\lambda_j$注水 & $\mathcal{N}(\bm{0},\bm{K}^{\star})$ \\
  \bottomrule
 \end{tabular}
\end{BookInlineTable}

这些结论都依赖指定的条件：噪声与输入独立，各次使用之间独立，噪声方差或协方差固定并为发送端所知，分量数固定，预算按每条码字计算。噪声在时间上相关、信道带有记忆或反馈时，需要随机过程的工具重新建立单次上界与大数律，可参考\BookChapterRef{chap:discrete-time-stochastic-processes}{卷九第3章“离散时间随机过程”}。容量定理是关于块长趋于无穷的陈述：图~\ref{fig:gcc-random-code}已经显示，有限块长下距离容量还有明显的差距，随机码本的存储和最近码字的搜索也随块长指数增长。下一章转向可以实际构造和译码的有限长度纠错码。
